%
%============= CHAPTER BREAK =================%
%
\chapter{Favorinus of Arles (ca 117--161 CE)}%
\vspace*{3mm} %
\justifying %

\noindent{}Favorinus flourished as a widely renowned Sophist and rhetor %
during the reign of Hadrian, even after the emperor publicly chastened him. %
Yet he outlived the emperor by over 20 years, %
and he continued to speak, write, and teach until his death. %
Intense scholarly debates continue to roil on about when Favorinus %
composed the oration \textit{On Exile} and %
whether it reflects authentic autobiography, %
caricatures his primary Sophist rival (Polemon of Laodicea|Smyrna), %
or is entirely fictive. %
Questions of autobiography and chronology go hand in hand; %
his relationship with Hadrian is another related factor. %
Here I dare not wade into the thick of such weeds, %
but simply adopt the position of past (Barigazzi) and current (Amato) Italian experts %
who affirm that this work of Favorinus is autobiographical in some sense, %
refers to his exile late in Hadrianʼs reign, %
and reflects a newfound confidence subsequent to Hadrianʼs death %
rather than a foolhardy attempt to take the emperor directly to task.\endnote{%
Barigazzi 1966; Amato 2005.} %

On this reading, Favorinusʼ reference to progeny, %
given his biological incapacity to father children, %
is not an authorial sticking point so much as a hang-up %
of modern scholars more preoccupied with functioning genitalia %
and patrilineage than Second Sophistic rhetorical conventions, %
patronage networks, and educational contexts. %
References to progeny and inheritance can be humorous, %
philosophical, communitarian, and/or adoptive in meaning, %
precisely the forms of fictive kinship that permeate the canonical Paulines. %
Aulus Gellius was among his beloved students %
and attested to Favorinusʼ ongoing, caring tutelage during %
the reign of Antoninus Pius. %
As we will soon see, there may also be neologistic grounds for a post-Hadrianic %
date for the composition. \clearpage %

\noindent{}\textbf{ἐπιπόθησις} / \textit{epipothēsis} / “yearning” as a lemma %
first appears in one of four roughly contemporaneous non-Christian texts, %
among whose candidates are Favorinusʼ oration \textit{On Exile}, % 
the \textit{Roman History} of Appian, %
and two grammatical works by Herodian. %
The corresponding verb, ἐπιποθέω / “to yearn”, goes back to Attic authors %
and was common in subsequent centuries (1649 TLG hits). %
The TLG indicates two previous, ambiguous hits for the verb or noun lemmas, %
but both are clearly verbs: %
μὴ ἔχων οὐκ ἐπιποθήσεις / “and not having you do not yearn” %
(Teles, \textit{Comparison of Poverty and Wealth});\endnote{%
Hense 1909:38.} %
οὐκ ἐπιποθήσεις ἐπ’ αὐτῷ οὐδ’ οὐ μὴ σκεπάσῃς αὐτόν / “you shall not %
yearn after him nor shall you cover him” (LXX Deut 13:9). %
For Favorinus, %
the noun (though separated by seven words from its definite article) %
supplies the syntactical opening of his speech: 

\vspace*{3mm} 
 \begin{quote}\textls[-20]{Ἡ δὲ | τῶν οἰκείων φίλων τε καὶ ξυγγενῶν | \textbf{ἐπιπόθησις} %
 τοῦ τῆς πατρίδος ἔρωτος ἐξηρ|τημένη δεύτερον ἐπ’ ἐκείνῳ ἀγώνισ|μα %
 προτέθειται.} \\[1mm] %
 \textls[-15]{Now the \textbf{yearning}, both for familiar friends and relatives, %
 surpassed by the lust for the fatherland, %
 is set forth as a second struggle after that.} %
 \end{quote} % 
 \vspace*{3mm} 

In the extant writings of Appian, %
the noun appears once, specifically in the volume entitled \textit{Celtica}, %
describing the Gallic sack of Rome (ca. 390 BCE) %
and written around the 150s CE. %
It is found within an excerpt preserved in the Byzantine \textit{Suda} (10th cent. CE), %
in a quotation derived in turn from the somewhat earlier \textit{Excerpta de Virtutibus}.\endnote{%
The \textit{Suda} entry is s.v., ὑφίσταται, and references \textit{Exc. de virt.} 9, p. 221.} %
In it, the Roman general Marcus Furius Camillus speaks %
in response to a plea by Caedicius, the legendary auditor of %
a divine oracle predicting the Gallic invasion. %

\vspace*{3mm} 
 \begin{quote}\textls[-25]{οὐκ ἂν εὐξάμην ἐπιποθῆσαί με Ῥωμαίους, %
 εἰ τοιαύτην ἤλπισα τὴν \textbf{ἐπιπόθησιν} αὐτοῖς ἔσεσθαι. %
 νῦν δὲ δικαιοτέραν εὐχὴν εὔχομαι, %
 γενέσθαι τῇ πατρίδι χρήσιμος ἐς τοσοῦτον ἀγαθοῦ, %
 ἐς ὅσον κακοῦ περιελήλυθεν.} \clearpage %
 I would not have prayed that the Romans yearn for me, %
 even if I hoped there would be such a \textbf{yearning} for them. %
 But now I pray a more righteous vow, %
 to become useful to fatherland for something of good, %
 to an extent equal to the evil it has endured. %
 \end{quote} % 
 \vspace*{3mm}

In his authentic, Antonine-era grammatical works, %
Herodian twice clarifies the lemmaʼs meaning, %
first briefly: ποθή ἡ \textbf{ἐπιπόθησις} / “longing is \textbf{yearning}” (\textit{orthog}. \textit{GG} 3.2:456). %
The second is more involved: %

\vspace*{3mm} 
 \begin{quote} ὥσπερ γὰρ χόλος τὸ πάθος, %
  χολὴ δὲ ἀφ’ οὗ τὸ πάθος, οὕτω καὶ πόθος τὸ πάθος, %
  ποθή δὲ ἡ \textbf{ἐπιπόθησις} καὶ ἐπιθυμία. (\textit{Pros.cath.} \textit{GG} 3.1:312)\\[1mm] %
  For just as anger (\textit{cholos}) is passion, %
  but bile (\textit{cholē}) what arises from passion, %
  so also longing (\textit{pothos}) is passion, %
  but longing (\textit{pothē}) is \textbf{yearning} or coveting. %
  \end{quote} % 
 \vspace*{3mm} %

The same noun lemma appears three times in the Pauline letters, %
but only in their canonical --- not Marcionite --- versions. %
The first two are related to the arrival of Titus in Corinth (2Cor 7:6--7, 11): %

\vspace*{3mm} 
 \begin{quote}\textls[-20]{ἀλλ’ ὁ παρακαλῶν τοὺς ταπεινοὺς παρεκάλεσεν %
 ἡμᾶς ὁ θεὸς ἐν τῇ παρουσίᾳ Τίτου· %
 ⁷ οὐ μόνον δὲ ἐν τῇ παρουσίᾳ αὐτοῦ ἀλλὰ καὶ
 ἐν τῇ παρακλήσει ᾗ παρεκλήθη ἐφ’ ὑμῖν, %
 ἀναγγέλλων ἡμῖν τὴν ὑμῶν \textbf{ἐπιπόθησιν}, τὸν ὑμῶν ὀδυρμόν, %
 τὸν ὑμῶν ζῆλον ὑπὲρ ἐμοῦ, ὥστε με μᾶλλον χαρῆναι. ... %
 ¹¹ ἰδοὺ γὰρ αὐτὸ τοῦτο τὸ κατὰ θεὸν λυπηθῆναι %
 πόσην κατειργάσατο ὑμῖν σπουδήν, %
 ἀλλὰ ἀπολογίαν, ἀλλὰ ἀγανάκτησιν, ἀλλὰ φόβον, %
 ἀλλὰ \textbf{ἐπιπόθησιν}, ἀλλὰ ζῆλον, ἀλλὰ ἐκδίκησιν·} \\[1mm] %
 But the God who encourages the humble encouraged us in the arrival of Titus. %
 ⁷ But not only in his arrival, but also in the %
 encouragement with which he was encouraged by us, %
 announcing to us the \textbf{yearning} of yours, %
 the lamentation of yours, %
 the rivalry of yours over me, so that I was gladdened even more. ... %
 ¹¹ For behold, this very thing, being pained based on God, %
 produced in you how much haste, %
 instead defense-speech, instead indignation, instead fear, %
 instead \textbf{yearning}, instead rivalry, instead vindication. \end{quote} % 
 \vspace*{1mm}

\noindent{}The third addresses Paulʼs planned journey to Spain (Rom 15:22--24): % 

 \vspace*{3mm} 
  \begin{quote}Διὸ καὶ ἐνεκοπτόμην τὰ πολλὰ τοῦ ἐλθεῖν πρὸς ὑμᾶς· %
    ²³ νυνὶ δὲ μηκέτι τόπον ἔχων ἐν τοῖς κλίμασι τούτοις, %
 \textbf{ἐπιποθίαν} δὲ ἔχων τοῦ ἐλθεῖν πρὸς ὑμᾶς ἀπὸ πολλῶν ἐτῶν, %
 ²⁴ ὡς ἂν πορεύωμαι εἰς τὴν Σπανίαν· \\[1mm] %
 \textls[-15]{Hence was I also hindered many times from coming unto you. %
  ²³ But now no longer having a place in these regions, %
  yet having a \textbf{yearning} for many years to come before you, %
  ²⁴ as I might go into Spain.}\end{quote} % 
  \vspace*{3mm}

Both Favorinus and Appian used the lemma with reference to patriotism, %
specifically the yearning for fatherland during exile or a time away from home. %
The usage in canonical Romans carries similar notes of geographical displacement, %
the longing for a place of familiar belonging, %
and perhaps even military-like duty to travel abroad. %
The two hits in canonical 2Cor %
bear the most resemblance to Favorinusʼ seminal usage. %
They (i.e., Favorinus) enjoyed a famously close association with Corinth, %
only for that fame to turn to public shame. %
It was in Corinth that an honorific statue of Favorinus was erected only later to be dethroned %
and replaced with another statue that mocked his physical appearance.\endnote{%
About this public scandal, %
see Favorinusʼ \textit{Corinthian Oration}, %
later wrongly attributed to Dio Chrysostom as \textit{Orat}. 37.} %
The use of this neologism in the canonical Paulines %
very likely casts Titus as a second Favorinus. %
In this regard, the notoriously ample, abnormal, and/or absent genitalia of Favorinus %
do at last turn out to be comparatively meaningful. %
For Titus was the associate whom Paul (Marcionite and canonical) refused to circumcise, %
thus the fictive Timothy of the canonical redaction had to undergo the procedure in his place. %